% ============================================================
%  thesis.tex — Main thesis file
%  Adapt this template to your institution's requirements.
%  See AGENTS.md for the workflow rules before editing.
% ============================================================

\documentclass[12pt, a4paper]{report}

% --- Packages ---
\usepackage[utf8]{inputenc}
\usepackage[T1]{fontenc}
\usepackage{geometry}
\geometry{margin=2.5cm}

\usepackage{setspace}
\onehalfspacing

\usepackage{graphicx}
\usepackage{booktabs}       % professional tables
\usepackage{hyperref}       % clickable references
\usepackage{cleveref}       % \cref{} instead of Figure~\ref{}
\usepackage{natbib}         % author-year citations; swap for biblatex if needed
\usepackage{amsmath, amssymb}
\usepackage{lipsum}         % placeholder text — remove before submission

% --- Metadata ---
\title{[Your Thesis Title]}
\author{[Your Name]}
\date{[Month Year]}

% ============================================================
\begin{document}

\maketitle
\tableofcontents
\listoffigures
\listoftables

% ============================================================
\begin{abstract}
[To be completed — write the abstract last.]
\end{abstract}

% ============================================================
\chapter{Introduction}
\label{ch:introduction}

% Problem statement
[To be completed — describe the problem and its significance.]

% Research questions
\section{Research Questions}
\label{sec:rq}

The main research question of this thesis is:

\begin{quote}
\textbf{RQ-Main:} [Your main research question]
\end{quote}

This is addressed through two sub-questions:

\begin{quote}
\textbf{RQ-1:} [First sub-question]\\
\textbf{RQ-2:} [Second sub-question]
\end{quote}

\section{Contributions}
\label{sec:contributions}

The main contributions of this thesis are:

\begin{enumerate}
    \item [Contribution 1]
    \item [Contribution 2]
    \item [Contribution 3]
\end{enumerate}

\section{Thesis Structure}
\label{sec:structure}

[To be completed — one sentence per chapter describing what each covers.]

\section{Delimitations}
\label{sec:delimitations}

[To be completed — what this thesis explicitly does not cover and why.]

% ============================================================
\chapter{Background and Related Work}
\label{ch:background}

\section{Technical Background}
\label{sec:background}

[To be completed]

\section{Related Work}
\label{sec:related}

[To be completed]

\section{Research Gaps}
\label{sec:gaps}

Based on the literature, three gaps motivate this thesis:

\begin{enumerate}
    \item \textbf{Gap 1} — [description] (aligned with RQ-1)
    \item \textbf{Gap 2} — [description] (aligned with RQ-2)
    \item \textbf{Gap 3} — [description] (aligned with RQ-Main)
\end{enumerate}

% ============================================================
\chapter{Methodology}
\label{ch:methodology}

\section{Research Design}
\label{sec:design}

[To be completed]

\section{Data}
\label{sec:data}

[To be completed]

\section{Pipeline}
\label{sec:pipeline}

[To be completed]

% Example figure placeholder
\begin{figure}[ht]
    \centering
    \fbox{\parbox{0.8\textwidth}{\centering\vspace{2cm}[Figure: pipeline overview — to be inserted]\vspace{2cm}}}
    \caption{Overview of the proposed pipeline.}
    \label{fig:pipeline}
\end{figure}

\section{Evaluation Metrics}
\label{sec:metrics}

[To be completed]

% ============================================================
\chapter{Results}
\label{ch:results}

% Results only: what was observed, not why.

\section{[Experiment 1 Name]}
\label{sec:results-exp1}

[To be completed]

% Example table placeholder
\begin{table}[ht]
    \centering
    \caption{[Experiment 1] results on [dataset].}
    \label{tab:results-exp1}
    \begin{tabular}{lcc}
        \toprule
        Method & Metric 1 & Metric 2 \\
        \midrule
        Baseline & TBD & TBD \\
        Proposed & TBD & TBD \\
        \bottomrule
    \end{tabular}
\end{table}

% ============================================================
\chapter{Discussion}
\label{ch:discussion}

% Discussion only: why / what it means. No new results here.

The results show [To be completed]. This suggests [To be completed] in relation to RQ-Main.

\section{Limitations}
\label{sec:limitations}

[To be completed]

% ============================================================
\chapter{Conclusion}
\label{ch:conclusion}

This thesis addressed [To be completed]. The main findings were [To be completed].

\section{Future Work}
\label{sec:future}

[To be completed]

% ============================================================
\bibliographystyle{apalike}     % change to your institution's required style
\bibliography{tex/biblio}

\end{document}
